\documentclass[11pt]{article}

\usepackage[margin=1in]{geometry}
\usepackage{amsmath,amssymb,amsfonts,amsthm}
\usepackage{booktabs}
\usepackage{graphicx}
\usepackage{hyperref}
\usepackage{xcolor}
\hypersetup{
    colorlinks=true,
    linkcolor=black,
    citecolor=black,
    urlcolor=blue!50!black,
    pdfborder={0 0 0},
    pdfborderstyle={/S/N},
}
\providecommand{\doi}[1]{\href{https://doi.org/#1}{doi:#1}}
\usepackage[expansion=false]{microtype}
\usepackage{authblk}
\usepackage[numbers,square,sort&compress]{natbib}

\newtheorem{theorem}{Theorem}
\newtheorem{proposition}[theorem]{Proposition}
\newtheorem{lemma}[theorem]{Lemma}
\newtheorem{corollary}[theorem]{Corollary}
\theoremstyle{definition}
\newtheorem{definition}[theorem]{Definition}
\newtheorem{assumption}[theorem]{Assumption}
\theoremstyle{remark}
\newtheorem{remark}[theorem]{Remark}
\newtheorem{observation}[theorem]{Observation}

\newcommand{\E}{\mathbb{E}}
\newcommand{\PP}{\mathbb{P}}
\newcommand{\Q}{\mathbb{Q}}
\newcommand{\R}{\mathbb{R}}
\newcommand{\C}{\mathcal{C}}
\newcommand{\X}{\mathcal{X}}
\newcommand{\KL}{D_{\mathrm{KL}}}
\newcommand{\GROUND}{\Gamma_{\mathrm{i}}}

\title{An Entropy-Regularized Growth-Optimal Selection Criterion \\
       for Discrete-Payoff Options: \\
       Theory and Empirical Validation on SP100 Credit Spreads}

\author{Peter J.~Mercurio, PhD}
\affil{\small \texttt{peter.mercurio@alumni.utoronto.ca}}
\date{\today}


\begin{document}
\maketitle

\begin{abstract}
The GROUND ratio is a portfolio-selection criterion that jointly
rewards Kelly growth and penalizes distributional divergence from a
uniform reference. In this work we introduce the intrinsic GROUND ratio by adopting the growth signal $g := \ln E$
(where $E$ is the Kelly expected value (EV)) to give the criterion a multiplicative
form that reads directly as a percent risk-adjusted return. Under this
choice, we observe that $\ln \GROUND = g - k \cdot \KL = J_k$ holds
exactly, where $J_k$ is the Hansen-Sargent multiplier-preferences functional. In the small-edge asymptotic limit where regularization weight $k \to 0$ and payoff edge shrinks to zero,
the reservation-price approximation under Kelly-optimal sizing is asymptotically proportional to
the agent's chosen belief-measure expectation, with proportionality constant $\tfrac{1}{2}$. As a special case, if that belief measure is
the minimum-entropy martingale measure (MEMM), the limiting approximation equals $\tfrac{1}{2}\E_{\Q^\star}[x]$. We apply the criterion to a
Greek-filtered weekly pool of Standard \& Poor's 100 vertical credit
spreads with $k = 20$ and selection threshold $\GROUND \geq 0.10\%$
frozen via in-sample grid sweeps before any 2025--2026 data is
scored. On the in-sample 2020--2024 window ($n=4{,}765$ candidates),
the strategy achieves Sharpe $1.85$ [Mertens 95\% CI on Sharpe: 1.33, 2.37].
On the strict 70-week 2025--2026 holdout, Sharpe $2.80$ [CI: 1.39, 4.21]. The extended
six-year sample yields Sharpe $2.07$ [CI: 1.56, 2.58]. Across
$6{,}130$ candidates, the top $\GROUND$ decile beats the bottom by
$32\times$ on mean per-trade yield, with the edge statistically
robust in holdout (Spearman $\rho$ rises from $+0.108$ in-sample to
$+0.153$ out-of-sample). The joint ranker outperforms both single-factor baselines (Kelly EV alone,
Kullback-Leibler divergence alone) on every metric, confirming that the intrinsic GROUND ratio is more powerful than either component alone.
\end{abstract}


\section{Introduction}\label{sec:intro}

The Kelly criterion (Kelly \cite{kelly1956}) prescribes the
growth-optimal sizing fraction for a repeatable wager under known
probabilities. When this criterion is applied to option-strategy
selection, two well-documented difficulties arise. The first is that
the probabilities are not actually known. They are estimated,
typically from implied-volatility-derived risk-neutral distributions
(Breeden and Litzenberger \cite{breeden1978}) or from realised
statistics, and small errors in those estimates can produce large
errors in the prescribed fraction (MacLean et al.\ \cite{maclean2011}). The second is that the Kelly objective rewards
implied distributions whose mass is sharply concentrated on a single
outcome. In a derivatives context this can mean rewarding implied
distributions whose sharpness reflects measurement noise as much as
genuine economic signal.

Both of these difficulties motivate a regularized version of the
Kelly objective, namely a criterion that retains the
growth-maximization core but penalizes distributional configurations
whose informational sharpness is suspect. A natural candidate for
the penalty is the Kullback-Leibler divergence of the implied
distribution from a maximally-uncertain reference, taken to be the
uniform distribution on the discrete outcome space. The GROUND
ratio of Mercurio et al.\ \cite{mercurio2020binary} took exactly this form, originally as
a comparative score: each candidate was ranked against a
reference candidate (typically the minimum-divergence candidate in
the pool) by a ratio of growth differences to KL differences.

The comparative form of the GROUND ratio is awkward in practice. The
selection of any one candidate depends on the reference candidate chosen, and ties at minimum have to
be handled procedurally. In this paper we replace the comparative
form with an intrinsic, per-candidate form. We then write the
intrinsic form in its canonical Kelly-EV-with-entropic-discount
expression (Section~\ref{sec:framework}), and apply it to a real
options market.

\paragraph{What the paper does.}
We make four contributions.

\begin{enumerate}
    \item We state the intrinsic GROUND ratio in its canonical
          Kelly-EV-with-entropic-discount form
          (Section~\ref{sec:framework}). Under the growth signal $g
          := \ln E$, where $E$ is the Kelly EV, the identity
          $\ln \GROUND = J_k$ holds exactly by construction. The
          multiplicative $\GROUND$ form and the additive Hansen-Sargent
          functional are therefore the same selection rule expressed
          in different units, and they produce identical selection on
          any menu of growth-positive candidates.
    \item We characterize the reservation-price approximation under
          the regularized log-growth objective. Under Kelly-optimal sizing,
          in the small-edge asymptotic limit the first-order approximation is
          proportional to the belief-measure expectation with constant $\tfrac{1}{2}$: $c^\star \approx \tfrac{1}{2} \E_p[x]$
          for any belief distribution $p$ (Theorem~\ref{thm:memm}). When $p$ is the
          minimum-entropy martingale measure of Frittelli \cite{frittelli2000}, the approximation equals
          $\tfrac{1}{2}\E_{\Q^\star}[x]$. This proportionality is a consequence of Kelly's growth-variance trade-off, not a derivation of MEMM.
    \item We validate the criterion on a real options market. To our
          knowledge this is the first end-to-end backtest of the
          intrinsic GROUND ratio on a real options market,
          including realistic candidate generation, filtering,
          sizing, and execution-cost treatment
          (Sections~\ref{sec:strategy} and~\ref{sec:results}). The
          ambiguity-aversion weight $k = 20$ and the threshold
          $\tau = 0.10\%$ were both selected on the in-sample
          window and frozen before any 2025--2026 data was scored;
          the canonical filter pipeline (DTE, delta, max-loss, OI,
          regime indicator, sizing, and probability split)
          was developed by informal exploration on the in-sample
          window and likewise frozen. See
          Section~\ref{sec:filters} for a full
          parameter-accounting paragraph.
    \item We isolate the contribution of each ranker component by
          replacing the intrinsic GROUND ratio with one factor at a
          time and re-running
          the strategy on the same candidate pool
          (Section~\ref{sec:comparison}). The full joint ranker
          outperforms both single-factor baselines on every metric
          we report. We take this as direct evidence that the joint
          ranker does work that neither component does on its own.
\end{enumerate}

\paragraph{Relation to our previous work.}
The use of information-theoretic quantities as risk measures in
portfolio optimization was proposed in Mercurio et al.\ \cite{mercurio2020portfolio},
where we argued that Shannon entropy is a natural alternative to
variance when the return distribution is discrete or otherwise
ill-suited to mean-variance treatment. We subsequently specialized
the construction to binary options in Mercurio et al.\ \cite{mercurio2020binary},
where we introduced the GROUND ratio as a portfolio-selection
criterion that simultaneously rewards Kelly growth and penalizes
deviation from a uniform reference distribution on the two-state
binary-wager simplex, and showed that the KL divergence of the
wager's outcome distribution against the uniform reference is a
convex risk measure in the sense of F\"ollmer and Schweizer
\cite{follmer2002}, with the uniform distribution playing the role
of the minimum-risk reference. The construction was subsequently
applied to credit-spread option portfolios in
Mercurio et al.\ \cite{mercurio2020}, which was the first use of an
entropy-based risk measure for portfolio optimization on credit
spreads under a generalized entropic framework.

The present paper builds directly on those three earlier papers
and extends them in three directions. First, we replace the
comparative GROUND form of Mercurio et al.\ \cite{mercurio2020binary}
with an intrinsic per-candidate form. By adopting the growth signal
$g := \ln E$, we obtain a multiplicative form that reads as a
percent risk-adjusted return, and we observe that this choice yields
an exact identity with the Hansen-Sargent multiplier-preferences
functional. Second, we characterize the reservation-price approximation under Kelly-optimal sizing.
In the small-edge asymptotic limit where $k \to 0$ and payoff edge shrinks, the
formula is asymptotically proportional to the belief-measure expectation with constant $\tfrac{1}{2}$
(Theorem~\ref{thm:memm}), a consequence of Kelly's growth-variance trade-off. As a special case, when $p$ is the minimum-entropy martingale
measure of Frittelli \cite{frittelli2000}, the limiting approximation equals $\tfrac{1}{2}\E_{\Q^\star}[x]$, not the full MEMM
expectation. This proportionality principle applies uniformly regardless of belief choice. Third, we apply the intrinsic GROUND ratio to SP100 weekly
vertical credit spreads and report a six-year end-to-end backtest,
including realistic candidate generation, Greek and liquidity
filtering, sizing, regime gating, and execution-cost treatment.

\paragraph{Relation to the entropic pricing literature.}
The Hansen-Sargent multiplier-preferences functional
Hansen and Sargent \cite{hansen2008robustness} and the variational preferences class
of Maccheroni et al.\ \cite{maccheroni2006ambiguity} provide the additive form $J_k =
g - k \KL$ to which the intrinsic GROUND ratio is exactly related by
exponentiation, $\GROUND = \exp(J_k)$. The incomplete-markets pricing literature
(Frittelli \cite{frittelli2000}; Delbaen et al.\ \cite{delbaen2002}; Fujiwara and Miyahara \cite{fujiwara2003}) studies
the minimum-entropy martingale measure as a canonical pricing kernel under distributional uncertainty.
Theorem~\ref{thm:memm} shows that GROUND, motivated from the Kelly growth optimization side,
exhibits an analogous entropy-regularization principle: under Kelly-optimal sizing, the reservation-price approximation is proportional
to any belief-measure expectation (proportionality constant $\tfrac{1}{2}$). This is not derived from MEMM,
but achieves a related robustness effect through Kelly's growth-variance trade-off. The intrinsic
GROUND ratio thus connects to robustness-preference theory (Hansen and Sargent \cite{hansen2008robustness}) and incomplete-markets pricing
(MEMM literature) through this proportionality principle, though from a different theoretical starting point.

\paragraph{Organization.}
The rest of the paper is organized as follows.
Section~\ref{sec:framework} states the canonical intrinsic
GROUND ratio $\GROUND$ and
the exact identity with the Hansen-Sargent functional.
Section~\ref{sec:reservation-theory} characterizes the reservation-price approximation
under Kelly-optimal sizing and its proportionality to belief-measure expectations.
Section~\ref{sec:strategy} describes the empirical setup: the SP100
universe, the candidate generation, the Greek filters, the sizing
rule, and the slippage overlay. Section~\ref{sec:results} reports
the headline results on the in-sample, holdout, and extended
samples, with Mertens confidence intervals.
Section~\ref{sec:ksweep} reports the in-sample $k$-sweep.
Section~\ref{sec:threshold} reports the in-sample threshold
sweep over $\tau$ and motivates the canonical
$\GROUND \geq 0.10\%$ selection rule.
Section~\ref{sec:comparison} reports the single-factor baseline
comparison. Section~\ref{sec:decile} reports ranker quality across
the full candidate distribution, split by sub-period.
Section~\ref{sec:discussion} discusses limitations and open
problems. Appendix~\ref{app:proofs} collects deferred proofs.


\section{The canonical intrinsic GROUND ratio}\label{sec:framework}

In this section we set up the framework and introduce the intrinsic
GROUND ratio $\GROUND$ in its canonical Kelly-EV-with-entropic-discount form.
We adopt a growth signal $g := \ln E$ to give the ratio a multiplicative
form that reads directly as a percent risk-adjusted return; under this choice,
the identity $\ln \GROUND = J_k$ holds by construction. We record the two
limiting cases of the ratio and the economic meaning of each.

\subsection{Setup}

Let $\X = \{\omega_1, \ldots, \omega_n\}$ be a finite outcome set
with $n \ge 2$, and let $\C = (x, p, c)$ be a contingent claim on
$\X$. Here $x: \X \to \R$ is the per-unit-stake payoff vector, $p$
is the probability vector on $\X$, and $c \ge 0$ is the
per-unit-stake cost of entry. For a wager fraction $w \in [0, 1)$
of the agent's bankroll, the post-bet bankroll multiplier in state
$\omega$ is $1 + w \, x(\omega)$, and the expected log-bankroll gain
across the outcomes is
\begin{equation}
\ell(w; p, x) \;:=\; \sum_{\omega \in \X} p(\omega) \, \ln\bigl(1 + w \, x(\omega)\bigr).
\label{eq:loggrowth}
\end{equation}
The Kelly-optimal fraction is $w^\star = \arg\max_w \ell(w; p, x)$;
existence and uniqueness on growth-positive claims are standard, and
we record them in Appendix~\ref{app:proofs}, Lemma~\ref{lem:wstar}.

At the optimum, two quantities are useful to keep distinct. The
first is the wealth multiplier per trade,
\begin{equation}
M(p, x) \;:=\; \exp\bigl(\ell(w^\star; p, x)\bigr) \;\ge\; 1.
\label{eq:M}
\end{equation}
The second is the Kelly EV, defined as the expected wealth
\emph{gain} per trade at Kelly-optimal sizing,
\begin{equation}
E(p, x) \;:=\; M(p, x) - 1 \;=\; \exp(\ell(w^\star)) - 1 \;\ge\; 0.
\label{eq:E}
\end{equation}
The Kelly EV is dimensionless, positive on growth-positive claims,
and zero on claims with $\E_p[x] \le 0$. We restrict the menu
throughout to growth-positive candidates ($E > 0$). Candidates with
$E \le 0$ have Kelly-optimal fraction $w^\star = 0$ and contribute
nothing.

It is worth pausing on the Kelly EV. Because $\ell(w^\star)$ is the
expectation of the concave function $\log(1 + w x)$, Jensen's
inequality gives $\ell(w^\star) \le \ln(1 + w^\star \E_p[x])$, with
strict inequality whenever $\mathrm{Var}_p[x] > 0$. Therefore $E$ is
bounded above by $w^\star \E_p[x]$ (the arithmetic Kelly EV) and is
strictly less whenever the payoff distribution has nonzero variance.
That is, the Kelly EV is variance-adjusted by construction; this is
what distinguishes it from the linear EV $\E_p[x]$. A high-EV but
high-variance trade has a lower Kelly EV than a moderate-EV but
low-variance trade with the same arithmetic mean.

We also need the relative entropy of $p$ from the uniform reference
$u_n(\omega) = 1/n$. In natural-log units (nats),
\begin{equation}
\KL(p \| u_n) \;=\; \sum_{\omega} p(\omega) \ln\!\frac{p(\omega)}{1/n} \;=\; \ln n - H(p),
\label{eq:KL}
\end{equation}
where $H(p)$ is the Shannon entropy. The divergence is in $[0, \ln
n]$. For the three-outcome credit-spread setting that occupies the
empirical portion of this paper, $n = 3$ and $\ln 3 \approx 1.099$.

\subsection{The intrinsic GROUND ratio}

\begin{definition}[Intrinsic GROUND ratio]\label{def:ground}
Fix $k > 0$ and a contingent claim $\C = (x, p, c)$ with $E(p, x) >
0$. The intrinsic GROUND ratio at parameter $k$ is
\begin{equation}
\boxed{\;
\GROUND(p, x, k) \;:=\; E(p, x) \cdot \exp\!\bigl(-k \cdot \KL(p \| u_n)\bigr).
\;}
\label{eq:ground}
\end{equation}
\end{definition}

The two factors carry the two roles. The first factor is the
variance-adjusted expected wealth gain per trade, that is, the
\emph{return} side of the ratio. The second factor $\exp(-k \cdot
\KL) \in (0, 1]$ is the entropic ambiguity discount, that is, the
\emph{risk} side of the ratio. The discount is $1$ when $p$
coincides with the uniform reference (zero divergence, full credit
to the Kelly EV), and decays exponentially as the implied
distribution becomes more concentrated. The product is dimensionless
and reads directly as the variance-adjusted expected return per trade
after the entropic ambiguity discount. For selected trades in our SP100 setting at the canonical
$k = 20$, the distribution of $\GROUND$ is right-skewed and
concentrated at small values: median $0.10\%$, mean $0.34\%$,
$90$th percentile $0.82\%$, $99$th percentile $4.0\%$, with a
small tail extending to roughly $13\%$. The exponential nature of
the entropic discount $\exp(-k \cdot \KL)$ at $k = 20$ is what
compresses most of the mass toward zero; only candidates whose
implied distribution sits very close to the uniform reference
escape the discount substantially.

The subscript $\mathrm{i}$ stands for ``intrinsic''. It distinguishes
the present form from the comparative form of Mercurio et al.\ \cite{mercurio2020binary},
in which a candidate $a$ was scored against a per-period reference
$b$ as described in the displayed equation below.
In the comparative form, a candidate $a$ was scored against a
per-period reference $b$ (taken to be the minimum-divergence
candidate in the pool) by
\begin{equation*}
\Gamma_{a \mid b} \;=\; \frac{\ell(w^\star_a;\, p_a, x_a) - \ell(w^\star_b;\, p_b, x_b)}{\KL(p_a \| u_n) - \KL(p_b \| u_n)},
\end{equation*}
which produces a per-pool score rather than a per-candidate one.
The intrinsic form is the evolution of that earlier construction
into a per-candidate, reference-free score, and it is the form we
adopt throughout the present paper.

\subsection{Exact identity with the Hansen-Sargent functional}

Define the growth signal as the log of the Kelly EV:
\begin{equation}
g(p, x) \;:=\; \ln E(p, x).
\label{eq:gdef}
\end{equation}
This is well-defined on growth-positive claims (where $E > 0$). With
this definition, $E = \exp(g)$, and the canonical
form~\eqref{eq:ground} becomes
\begin{equation*}
\GROUND(p, x, k) \;=\; \exp(g(p, x)) \cdot \exp\bigl(-k \cdot \KL(p \| u_n)\bigr).
\end{equation*}
Taking the natural logarithm of both sides yields the additive form:
\begin{equation}
\ln \GROUND(p, x, k) \;=\; g(p, x) - k \cdot \KL(p \| u_n) \;=:\; J_k(p, x),
\label{eq:Jk}
\end{equation}
which is the multiplier-preferences functional of
Hansen and Sargent \cite{hansen2008robustness} and the variational preferences class
of Maccheroni et al.\ \cite{maccheroni2006ambiguity}. Equivalently, in multiplicative
form,
\begin{equation}
\GROUND(p, x, k) \;=\; \exp\bigl(J_k(p, x)\bigr).
\label{eq:exact-identity}
\end{equation}

This identity holds exactly. It is an algebraic consequence of the
choice $g := \ln E$, not an approximation, and the derivation
required no Taylor expansion. The growth signal $g$ defined here is
not the classical Kelly log-growth $\ell(w^\star)$; the two are
related by $g = \ln(\exp(\ell(w^\star)) - 1)$ but differ
substantively in their typical numerical range. Classical Kelly
log-growth $\ell(w^\star)$ is a small positive number (roughly
$0.01$ to $0.10$ in our setting), while $g = \ln E$ is a large
negative number (roughly $-3$ to $-7$ in nats). The exponentiation
in $\GROUND = \exp(J_k)$ undoes the log, returning a small positive
risk-adjusted EV.

\paragraph{Why $g := \ln E$ and not $\ell(w^\star)$.}
The choice $g := \ln E$ is not the unique way to write an additive
form for $\GROUND$. A natural alternative is to use the classical
Kelly log-growth $\ell(w^\star)$ as the growth signal, in which
case the additive form becomes $\ell(w^\star) - k \cdot \KL$, which
also coincides with Hansen-Sargent multiplier preferences under a
different identification. We chose $g := \ln E$ for one reason: it
is the only choice under which the \emph{multiplicative}
$\GROUND$, which is what a trader actually reads as a number on a
screen, is interpretable as a percent risk-adjusted expected
return per trade. The Kelly EV $E = \exp(\ell(w^\star)) - 1$ is
the per-trade wealth gain in proportional units; multiplying it by
the entropic discount $\exp(-k \cdot \KL)$ gives a quantity in the
same units. Under the alternative $g = \ell(w^\star)$, the
displayed score would be the wealth multiplier $M = \exp(\ell)$
times the entropic discount, which is not a percent return and
requires the reader to subtract $1$ mentally. The exact
Hansen-Sargent identity in \eqref{eq:exact-identity} is a
consequence of the economically motivated choice, not the
motivation itself.

\begin{observation}[Selection equivalence]\label{obs:selection}
On any menu of growth-positive claims, ranking by the multiplicative
form~\eqref{eq:ground} returns the same candidate as ranking by the
additive form~\eqref{eq:Jk}. That is, $\arg\max_i \GROUND(p_i, x_i,
k) = \arg\max_i J_k(p_i, x_i)$.
\end{observation}

The observation is a one-line consequence of the identity
$\GROUND = \exp(J_k)$ composed with strict monotonicity of the
exponential. We record it explicitly because the two forms are
useful in different contexts. The multiplicative form is convenient
empirically, because the displayed value reads as a per-trade
percent return. The additive form is convenient analytically,
because the limits in $k$ are cleaner in log-space, and the pricing
derivation in Section~\ref{sec:reservation-theory} uses the additive form
directly.

\subsection{Limiting cases}\label{sec:limits}

Two limits of the canonical ratio are immediate.

\begin{proposition}[Kelly limit]\label{prop:kelly-limit}
On any finite menu of growth-positive claims, as $k \to 0^+$ the
$\GROUND$-optimal selection $\arg\max_i \GROUND(p_i, x_i, k)$ converges
to the pure Kelly-EV-optimal selection $\arg\max_i E(p_i, x_i)$.
\end{proposition}

\begin{proposition}[Maximum-entropy limit]\label{prop:maxent-limit}
On any finite menu of growth-positive claims, as $k \to \infty$ the
$\GROUND$-optimal selection converges to the minimum-divergence claim,
$\arg\min_i \KL(p_i \| u_n)$.
\end{proposition}

Both follow from continuity together with the boundedness of $E$
and $\KL$ on a finite menu; proofs are in
Appendix~\ref{app:proofs}. Between the two limits the ratio
interpolates monotonically in $k$. At $k = 0$ the agent ignores
distributional concentration and selects on pure Kelly EV; at
$k \to \infty$ the agent ignores growth and selects the most uniform
implied distribution; for intermediate $k$, the agent trades the
two off against each other. The empirical $k$-sweep
(Section~\ref{sec:ksweep}) explores this interpolation on the
in-sample window.


\section{Reservation prices under regularized growth}\label{sec:reservation-theory}

Section~\ref{sec:framework} addresses selection: which claims should an agent ranking by
GROUND choose from a menu? We now address the complementary pricing question: what
is the agent's reservation price (breakeven cost) for a single claim under the
regularized log-growth objective?

The regularized reservation price balances two forces. The Kelly log-growth $\ell(w^\star)$
measures expected wealth gain. The entropy penalty $k \cdot \KL(p \| u_n)$ discounts prices
proportional to how sharply the belief $p$ departs from the uniform reference. In the
small-edge asymptotic limit where payoff edge shrinks and $k \to 0$, the first-order
approximation exhibits a universal proportionality principle: the reservation price scales
with any belief-measure expectation with proportionality constant determined by Kelly's
growth-variance trade-off. This proportionality property connects GROUND to classical
utility-based pricing and the broader robustness-preference literature (Hansen and Sargent \cite{hansen2008robustness}).

\subsection{Reservation price under regularized growth}

For the reservation price we use the classical Kelly log-growth
$\ell(w^\star; p, x)$ from~\eqref{eq:loggrowth}, rather than the
growth signal $g = \ln E$ used for selection. The utility
indifference price is the standard log-utility object, and
$\ell$ admits the small-bet Taylor expansion that we will need in
the proof below.

\begin{definition}[$k$-regularized reservation price]\label{def:reservation}
Fix $k > 0$ and a contingent claim $\C = (x, p, c)$ with
$\ell(w^\star; p, x) - k \cdot \KL(p \| u_n) > 0$. The $k$-regularized reservation price is
the highest cost $c^\star_k$ at which the agent's regularized
log-utility, entered at the Kelly-optimal fraction $w^\star(p,x)$ computed
on the unregularized claim, equals the no-trade baseline:
\begin{equation}
c^\star_k(p, x) \;=\; \sup \Bigl\{ c \ge 0 \;:\; \ell(w^\star; p, x) \;-\; c \cdot w^\star(p, x)
    \;-\; k \cdot \KL(p \| u_n) \;\ge\; 0 \Bigr\}.
\label{eq:reservation}
\end{equation}
Solving at equality gives the closed form
\begin{equation}
c^\star_k(p, x) \;=\; \max\left(0, \frac{\ell(w^\star; p, x) - k \cdot \KL(p \| u_n)}{w^\star(p, x)}\right).
\label{eq:reservation-form}
\end{equation}
\end{definition}

The denominator $w^\star$ converts the per-bankroll growth quantity
in the numerator into a per-unit-stake price. Note that $w^\star$ is held
fixed (pre-cost), not re-optimized as cost $c$ varies; this is a first-order approximation,
not exact utility-indifference pricing. In the unregularized
case $k = 0$, the formula reduces to $c^\star_0 = \ell(w^\star) / w^\star$, a first-order log-utility reservation-price approximation (cf.\ MacLean et al.\ \cite{maclean2011}), with the same limitation of holding $w^\star$ fixed pre-cost.

\subsection{Asymptotic convergence in the small-edge limit}

The intrinsic GROUND-induced reservation-price formula~\eqref{eq:reservation-form} is a first-order
approximation to exact utility-indifference pricing. Under Kelly-optimal sizing, when the payoff edge shrinks
and regularization weight vanishes, the formula becomes asymptotically proportional to the expected payoff
under the agent's belief distribution, with proportionality constant $\tfrac{1}{2}$ (reflecting the growth-variance trade-off).
The next theorem makes this precise for any belief $p$, including the special case where $p$ is the minimum-entropy martingale
measure $\Q^\star$ of Frittelli \cite{frittelli2000}.

\begin{theorem}[Intrinsic GROUND-induced reservation price: first-order approximation in the small-edge limit]\label{thm:memm}
Let $\PP = u_n$ be the agent's reference probability measure.
Consider a sequence of claims $\C_m = (x_m, p_m, c_m)$ where payoffs $x_m$ are uniformly bounded,
edges $\mu_m := \E_{p_m}[x_m] \to 0^+$, and second moments are uniformly bounded away from zero,
$0 < \sigma_{\min}^2 \le \E_{p_m}[x_m^2] \le \sigma_{\max}^2 < \infty$. Under Kelly-optimal sizing
$w^\star_m \approx \mu_m / \E_{p_m}[x_m^2]$, if $k_m \to 0^+$ and $k_m = o(\mu_m^2)$, then
\begin{equation*}
c^\star_{k_m}(p_m, x_m) \;=\; \tfrac{1}{2}\,\E_{p_m}[x_m] \;+\; o(\E_{p_m}[x_m]).
\end{equation*}
Thus the first-order reservation-price approximation is proportional to the belief-measure expectation, with proportionality constant $\tfrac{1}{2}$ under Kelly-optimal sizing.
When $p$ is the minimum-entropy martingale measure $\Q^\star$ of Frittelli \cite{frittelli2000},
the limiting approximation equals $\tfrac{1}{2}\E_{\Q^\star}[x]$, not the full MEMM expectation. This proportionality
reflects Kelly's growth-variance trade-off, not a special relationship to MEMM: any belief measure
yields the same proportionality constant under Kelly-optimal sizing.
\end{theorem}

\begin{proof}
We work on the finite simplex $\Delta_n$. The formula~\eqref{eq:reservation-form} is a
first-order Taylor approximation to exact utility-indifference pricing: cost is subtracted linearly
outside the log rather than inside the payoff/bankroll multiplier, and $w^\star$ is fixed
(pre-cost) rather than re-optimized. Despite these approximations, the limit holds.

From~\eqref{eq:reservation-form},
\begin{equation*}
c^\star_k(p, x) \;=\; \frac{\ell(w^\star; p, x) - k \, \KL(p \| u_n)}{w^\star(p, x)}.
\end{equation*}
Under Kelly-optimal sizing $w^\star_m \approx \mu_m / E[x_m^2]$, Taylor-expand $\ell(w; p, x) = \E_p[\ln(1 + wx)]$:
\begin{equation*}
\ell(w^\star; p, x) \;=\; w^\star \mu - \tfrac{(w^\star)^2}{2} E[x^2] + O((w^\star)^3).
\end{equation*}
Divide by $w^\star$:
\begin{equation*}
\frac{\ell(w^\star)}{w^\star} \;=\; \mu - \tfrac{w^\star}{2} E[x^2] + O((w^\star)^2) \;=\; \mu - \tfrac{\mu}{2} + O(\mu^2) \;=\; \tfrac{\mu}{2} + O(\mu^2).
\end{equation*}

The entropy penalty is $k \cdot \KL(p \| u_n) / w^\star \approx k \cdot E[x^2] \cdot \KL / \mu$.
When $k = o(\mu^2)$, this term is $o(\mu)$ and negligible relative to the $\mu/2$ leading term. Thus:
\begin{equation*}
c^\star_k = \frac{\ell(w^\star) - k \cdot \KL}{w^\star} \;=\; \tfrac{\mu}{2} + O(\mu^2) - o(\mu) \;=\; \tfrac{1}{2}\E_p[x] + o(\E_p[x]).
\end{equation*}
\qed
\end{proof}

\begin{remark}
The key insight of Theorem~\ref{thm:memm} is the universal proportionality principle: under Kelly-optimal sizing,
the first-order reservation-price approximation scales with any belief-measure expectation with proportionality constant 1/2.
This constant reflects Kelly's fundamental growth-variance trade-off (Section~\ref{sec:framework}).
The proportionality holds uniformly regardless of whether the belief $p$ is risk-neutral, empirical, or the minimum-entropy martingale measure.

For $k > 0$ (the practical regime), the entropy penalty $k \cdot \KL(p \| u_n)$ further discounts prices in proportion to
distributional divergence from uniform. This is a robustness mechanism: it down-weights beliefs
that are sharply concentrated due to estimation noise or model misspecification. The entropy penalty
strength $k$ is chosen once in-sample (Section~\ref{sec:ksweep}) and applies identically to any belief measure.

\emph{Historical note:} Incomplete-markets pricing literature (Frittelli \cite{frittelli2000}, Delbaen et al.\ \cite{delbaen2002})
studies the minimum-entropy martingale measure as a canonical pricing kernel. The proportionality principle
shows that GROUND applies an analogous entropy regularization, but agnostically with respect to belief—not derived from MEMM, but
achieving a similar distributional robustness effect through a different mechanism (Kelly's growth optimization).
\end{remark}


\section{Empirical setup}\label{sec:strategy}

The rest of the paper validates the criterion empirically on a real
options market. This section describes the data, the candidate
generation, the filters, the sizing rule, and the slippage
treatment, all of which are held constant across every reported
result. A reader who is principally interested in the headline
results may skip directly to Section~\ref{sec:results}.

\subsection{The universe and the outcome space}

We backtest on the SP100 universe over the period 2020-01-01
through 2026-05-04. The data is the standard daily option chain,
recording bid, ask, mid, open interest, and implied volatility for
each contract. Each Monday we construct a candidate menu of
vertical credit spreads (bull put and bear call) on each
underlying. Each spread has three terminal outcomes at expiry:
\begin{itemize}
    \item \emph{full win}, where the underlying expires beyond the
        short strike in the favourable direction, and the trader
        keeps the entire credit;
    \item \emph{full loss}, where the underlying expires beyond
        the long strike on the adverse side, and the trader loses
        the maximum loss;
    \item \emph{partial}, where the underlying expires in the zone
        between the two strikes, and the trader keeps some
        fraction of the credit.
\end{itemize}
The outcome space is therefore $\X = \{+b,\, +\alpha(b) b,\, -1\}$
in per-dollar-of-risk units, where $b = \text{credit} /
\text{max-loss}$ and $\alpha(b) = (b - 1) / (2b)$ is the principled
partial multiplier (the uniform-in-strikes linear partial average
between $-1$ and $+b$). The three-outcome structure is what
distinguishes the credit-spread setting from the two-outcome binary
wagers of Mercurio et al.\ \cite{mercurio2020binary}.

The universe is the SP100.

\subsection{Greek filters and sizing}\label{sec:filters}

The canonical filter set is held constant across every reported
result. The constraints are:
\begin{itemize}
    \item days to expiry (DTE) in $[3, 8]$ trading days;
    \item short-leg absolute delta in $[0.35, 0.65]$, the
        moderate-moneyness sweet spot;
    \item per-spread maximum loss not exceeding $\$5$ per share;
    \item both legs at open interest (OI) of at least $100$
        contracts, a basic liquidity threshold;
    \item a binary regime gate on the SPDR S\&P 500 ETF Trust
        (ticker symbol SPY). SPY's close must be above its 100-day
        simple moving average (SMA) for a bull-put candidate to be
        eligible, and below the same SMA for a bear-call candidate.
\end{itemize}
This is the only macro filter in the canonical configuration.
Other filters that we tried, including per-ticker regime gates,
gap filters, filters conditioned on the CBOE Volatility Index (VIX),
theta-to-credit ratio filters, and credit-ratio caps, were either
redundant with the canonical set or return-degrading. We keep only
the canonical set.

Sizing is flat: every selected trade is entered at $q = 1$ or
$q = 2$ contracts per spread (the two reported configurations).
A per-trade ceiling of $\texttt{MAX\_CONTRACTS} = 50$ guards
against pathological cases under fractional-Kelly sizing variants
that are evaluated in robustness checks but not used in the
canonical configuration; with flat $q = 1$ or $q = 2$ the ceiling
never binds. The Kelly-optimal fraction $w^\star$ is used only as
a ranking input and never to size positions.

Each week, the strategy selects every candidate from the surviving
pool whose $\GROUND$ score clears a fixed cutoff
$\tau = 0.10\%$. The cutoff has a clean economic reading because
$\GROUND$ is dimensionless: setting $\GROUND \geq \tau$ is
equivalent to imposing a $\tau$-per-trade hurdle on the
variance-adjusted expected return after the entropic ambiguity
discount. We adopt $\tau = 0.10\%$ as the canonical threshold;
Section~\ref{sec:threshold} reports the in-sample threshold sweep
that motivates the choice. The threshold selects approximately
$2.8$ trades per week from the surviving pool on the in-sample
window.

\subsection{Probability estimation}

The per-leg in-the-money (ITM) probabilities are recovered from
each leg's delta via the standard delta-as-probability
approximation
\begin{equation*}
\Pr(\text{ITM at expiry}) \;\approx\; |\Delta|.
\end{equation*}
For short-dated options this diverges from the true physical
probability under any non-zero drift, and the implied-volatility
smile introduces additional moneyness-dependent error. In earlier
iterations of this work we tested realised-volatility,
drift-adjusted, and blended IV-RV estimators as alternatives.
None produced a meaningful improvement in Sharpe under the
$\GROUND$ ranker, and we revert to the delta-based estimator for
its simplicity and its independence from a vol-history
calibration window.

The full-win, partial, and full-loss probabilities
$(p, r_o, q)$ on the three-outcome space follow directly from the
two legs' deltas. For a bull-put spread with short strike
$K_s$ and long strike $K_\ell < K_s$, both legs are puts and the
ITM probabilities are nested: $\Pr(S_T < K_\ell)$ is the
full-loss event (the underlying expires past the long strike),
$\Pr(S_T < K_s)$ is the any-loss event, and their difference is
the partial-loss mass. Under the delta-as-probability
approximation this gives
\begin{equation*}
q \;=\; |\Delta_\ell|, \qquad
r_o \;=\; |\Delta_s| - |\Delta_\ell|, \qquad
p \;=\; 1 - |\Delta_s|,
\end{equation*}
where $\Delta_s$ and $\Delta_\ell$ are the short- and long-leg
deltas respectively. The bear-call construction is symmetric
with call deltas. No free splitting parameter is involved; the
broker-quoted long-leg delta directly supplies the partial-vs-full
loss decomposition, with the market's strike-specific IV
calibration baked in.

These probabilities feed $\ell(w^\star)$ and $\KL(p \| u_3)$ for
each candidate, both computed in natural-log units (nats).

\paragraph{Free-parameter accounting.}
For transparency, the canonical configuration has two parameters
selected from data on the in-sample window via documented sweeps
and frozen before any holdout scoring: the ambiguity-aversion
weight $k = 20$ (Section~\ref{sec:ksweep}, nine-point sweep) and
the threshold $\tau = 0.10\%$ (Section~\ref{sec:threshold},
eight-point sweep). The remaining numerical settings,
DTE $\in [3, 8]$, short-leg $|\Delta| \in [0.35, 0.65]$, max-loss
$\le \$5$/share, OI $\ge 100$, the SPY 100-day SMA regime indicator,
and the flat sizing rule, are filter and
configuration choices made during development. We do not claim these were never adjusted on the
in-sample window; alternative filters were tried (per-ticker
regime gates, gap filters, VIX-conditional filters,
theta-to-credit ratio filters) and found either redundant or
return-degrading, and the surviving filter set was retained for
the canonical configuration. The entire filter pipeline should
therefore be regarded as in-sample-selected via informal
exploration, with $k$ and $\tau$ as the only quantities selected
via formal grid sweeps. The strict 2025--2026 holdout was
scored with the entire pipeline frozen.

\subsection{Execution model and slippage}\label{sec:slippage}

The headline profit and loss (P\&L) is computed at the bid-ask
mid-price for both legs, held to expiry. To stress-test the
result against execution costs, we also report a post-hoc
slippage overlay of $3$ cents per leg. The vertical credit spread
has two legs, so the per-share slippage cost at $3$ cents per leg
is $6$ cents per share. At the standard $100$-share multiplier
this works out to $\$6$ per spread for the $q = 1$ configuration
and $\$12$ per spread for the $q = 2$ configuration. The overlay
is one-sided: spreads are held to expiry and settle at intrinsic
value, so there is no closing fill and no exit slippage. We
intend the $3$-cent figure as a conservative upper bound for an
investor capable of working a limit-order ladder in a
normally-quoted SP100 option chain. Realised per-leg slippage on
SP100 weekly options worked through a limit-order ladder is
typically below $1$ cent; the $3$-cent overlay is a worst-case
ceiling intended to bound transaction-cost risk in the headline
reporting rather than to estimate expected execution costs.
Selection is conducted at the mid-price; the slippage overlay is
applied to realised P\&L and to the wagered denominator only,
since the haircut is part of the real dollar at risk on entry.


\section{Headline results}\label{sec:results}

This section reports the strategy's performance on the three time
windows of interest: the in-sample window on which $k$ was
selected, the strict out-of-sample holdout window, and the
extended sample that combines the two. We then report Mertens 95\%
confidence intervals on the Sharpe ratios, compare against an SPY
buy-and-hold benchmark, and quantify the impact of the slippage
overlay.

\subsection{In-sample, holdout, and extended sample}

The in-sample 2020--2024 window is the window on which $k = 20$
was selected from the sweep over $\{0.5, 1, 2, 5, 10, 20, 50,
100, 200\}$ described in Section~\ref{sec:ksweep}. The holdout
window covers 2025-01-01 through 2026-05-04 over 70 weeks, and was
never used in $k$-selection. The extended sample is the union of
the two windows, and it is the largest-$T$ row in the headline
table.

\begin{table}[htbp]
\centering
\small
\begin{tabular}{lrrrrrr}
\toprule
window           & Trades & Final & Sharpe & Calmar & Max DD & Yield (0\textcent) \\
\midrule
\textbf{$q = 1$ (flat, 1 contract per spread)} \\
in-sample 2020--2024   & 731   & \$24{,}653 & 1.85 & 19.03 & $-7.7\%$  & 29.21\%  \\
holdout 2025--2026     & 215   & \$17{,}116 & 2.80 & 14.23 & $-5.0\%$  & 35.55\%  \\
extended 2020--2026    & 946   & \$31{,}769 & 2.07 & 16.63 & $-7.7\%$  & 31.08\%  \\
\midrule
\textbf{$q = 2$ (flat, 2 contracts per spread)} \\
in-sample 2020--2024   & 731   & \$39{,}306 & 1.85 & 19.03 & $-15.4\%$ & 29.21\%  \\
holdout 2025--2026     & 215   & \$24{,}232 & 2.80 & 14.23 & $-10.0\%$ & 35.55\%  \\
extended 2020--2026    & 946   & \$52{,}538 & 2.07 & 16.63 & $-15.4\%$ & 31.08\%  \\
\bottomrule
\end{tabular}
\caption{Headline statistics for $\GROUND$ at $k = 20$ with the
canonical $\GROUND \geq 0.10\%$ threshold selection rule on SP100
weekly credit spreads, at both flat sizing levels. The holdout row
is the strict out-of-sample evaluation: both $k = 20$ and
$\tau = 0.10\%$ were frozen from the in-sample window before any
holdout data was scored. Yields are sportsbook-style
$\sum \mathrm{P\&L} / \sum \mathrm{wagered}$ ratios. The Sharpe
ratio is annualised at $\sqrt{52}$ on weekly returns scaled to a
$\$10{,}000$ starting bankroll. Max DD is the peak-to-trough
drawdown on bankroll. The $q = 2$ rows are the $q = 1$ rows with
doubled trade size; P\&L, drawdown, and final equity scale
linearly, while Sharpe, yield, and Calmar are sizing-invariant by
construction. The high Calmar ratios reflect the strategy's small peak-to-trough
drawdowns (7.7\% in-sample, 5.0\% holdout) relative to annualised returns;
under Kelly-coupled compounding (rather than flat sizing), the Calmar would moderate accordingly.}
\label{tab:headline}
\end{table}

\subsection{Holdout compared with in-sample}

On the strict 2025--2026 holdout, the strategy posts a Sharpe of
$2.80$ and a yield of $35.55\%$ at zero slippage. Both are
\emph{higher} than the in-sample figures of $1.85$ and $29.21\%$,
which is the opposite of what one usually finds. The out-of-sample outperformance reflects the uncapped b-credit configuration adopted in the current version: high-b trades that would have been capped in earlier configurations enter the holdout set, and while they experience higher loss frequency, their occasional large wins drive overall wealth higher. Caution is warranted
before over-interpreting this out-of-sample out-performance. The holdout is short (70 weeks), and the
2025--2026 window happened to be a favourable environment for the
credit-spread direction selected by the SPY trend filter (SPY
closed above its 100-day SMA for most of the holdout). We report
Mertens confidence intervals on the holdout Sharpe in
Section~\ref{sec:sharpe-ci}, which give a more cautious read on
how tight the holdout estimate really is.

\paragraph{The April 2025 tariff-shock window.}
The 70-week holdout includes the late-March and April 2025
tariff-shock selloff, the largest single market disruption of the
period. We note two facts about how the strategy fared. First, the
single largest weekly loss of the entire holdout was $-2.38\%$, on
the week of 2025-03-17 (the entry into the tariff window). Set
against the in-sample single-week loss distribution, this is substantially
less severe than the in-sample worst-week loss of $-6.59\%$ (the week of
2020-04-13, in the COVID drawdown), and falls near the in-sample
$1$st-percentile loss of $-2.65\%$ (the holdout worst week sits at
the $\sim 2.3$rd percentile of in-sample weekly returns). The
tariff-window worst week is, in other words, an unremarkable
sub-2.4\% draw by the standards of what the in-sample window
already revealed. Second, the direction-conditional regime gate,
which prevented bear-call entries during the post-shock SPY
downtrend, automatically cut exposure during the most volatile
bear days. Whether one calls this a qualitative consistency check
or a structural feature of the strategy, the tariff window did
not break it.

\subsection{Confidence intervals on Sharpe}\label{sec:sharpe-ci}

Sharpe ratios estimated from short samples have wide confidence
intervals, particularly under non-Gaussian return distributions.
Lo \cite{lo2002sharpe} gave a closed-form standard error under an
iid Gaussian-returns assumption, and Mertens \cite{mertens2002sharpe}
extended that formula to handle arbitrary skewness and excess
kurtosis under iid returns. The Mertens standard error depends on
the third and fourth standardised moments of the return
distribution: skewness ($\gamma_3$, asymmetry of the
distribution; positive values indicate a right tail) and excess
kurtosis ($\gamma_4$, fat-tailedness relative to the normal
distribution; positive values indicate fatter tails). The
strategy's weekly returns are right-skewed (occasional large wins
mixed with frequent small losses) and fat-tailed, so Mertens
intervals are the appropriate primary reference. We report
Mertens 95\% intervals on all three windows in
Table~\ref{tab:sharpe-ci}.

\begin{table}[htbp]
\centering
\small
\begin{tabular}{lrrrrcc}
\toprule
window                & $T$ & Sharpe & Skew $\gamma_3$ & Excess kurt $\gamma_4$ & $\sigma_{\hat S}$(Mertens) & 95\% CI \\
\midrule
in-sample 2020--2024  & 260 & 1.85 & $+1.72$ & $+6.95$ & 0.27 & $[1.33,\,2.37]$ \\
holdout 2025--2026    &  70 & 2.80 & $+1.93$ & $+7.42$ & 0.72 & $[1.39,\,4.21]$ \\
extended 2020--2026   & 330 & 2.07 & $+1.80$ & $+7.23$ & 0.26 & $[1.56,\,2.58]$ \\
\bottomrule
\end{tabular}
\caption{Annualised Sharpe with iid-Gaussian
\cite{lo2002sharpe} standard errors extended to arbitrary higher
moments via Mertens \cite{mertens2002sharpe}. $T$ is the number
of Monday-entry weeks; $\gamma_3$ is sample skewness; $\gamma_4$
is sample excess kurtosis (full kurtosis minus 3, so a normal
distribution has $\gamma_4 = 0$). Sharpe is annualised at
$\sqrt{52}$.}
\label{tab:sharpe-ci}
\end{table}

Three observations follow from the intervals.

First, the holdout's Mertens interval $[1.39, 4.21]$ excludes zero
by a wide margin. Even its lower edge ($1.39$) sits above the SPY
matched-window point Sharpe of $1.02$ from Table~\ref{tab:spy}. We
do not claim the result is definitive on 70 weeks, but it is not
consistent with luck under the iid assumption.

Second, the in-sample interval $[1.33, 2.37]$ is tighter, but it
carries the usual caveat that both $k$ and the threshold $\tau$
were selected on this window. Therefore, the in-sample Sharpe is
not a clean estimate of the strategy's true Sharpe; it is an
upper-biased estimate.

Third, the extended-sample row is the most defensible single
number for external communication. It has the largest $T$ ($330$
weekly observations), the tightest confidence interval, and it
combines the parameter-selected window with the strict holdout.
Its lower 95\% bound is $1.56$, which is above SPY's same-period
Sharpe across the entire reported range. For reporting outside
the paper, we use this row.

\subsection{SPY benchmark}

The natural benchmark is SPY buy-and-hold over the matched
windows. We report it in Table~\ref{tab:spy}.

\begin{table}[htbp]
\centering
\small
\begin{tabular}{lrrr}
\toprule
window                   & Sharpe & Total ROI & Max DD \\
\midrule
SPY 2020--2024           & 0.77 & $+94.6\%$  & $-31.8\%$ \\
SPY 2025--2026 (matched) & 1.02 & $+24.6\%$  & $-16.9\%$ \\
SPY 2020--2026 (matched) & 0.82 & $+141.8\%$ & $-31.8\%$ \\
\bottomrule
\end{tabular}
\caption{SPY buy-and-hold benchmark over the matched windows.
Return on investment (ROI) is the total period return. Sharpe is
annualised at $\sqrt{52}$. See the convention paragraph below.}
\label{tab:spy}
\end{table}

\paragraph{SPY benchmark convention.}
For reproducibility we pin the SPY benchmark convention used in
this paper. The price series is SPY daily close, dividend-adjusted
at the source (back-adjusted for distributions through the
end-of-series date). Weekly returns are computed from the
W-Friday resampled close (last Friday close at or before the
period boundary). Window boundaries are inclusive on both ends:
the in-sample window runs from the first W-Friday on or after
2020-01-01 through the last W-Friday on or before 2024-12-31; the
holdout runs from the first W-Friday on or after 2025-01-01
through the last W-Friday on or before 2026-05-04; the extended
window is their union. The Sharpe and ROI figures of
Table~\ref{tab:spy} are the canonical numbers under this
convention and should be reproducible by any reader using the same
W-Friday/adjusted-close pipeline.

The strategy's Sharpe exceeds the SPY benchmark's on every window,
with a substantial margin: $1.85$ versus $0.77$ in-sample, $2.80$
versus $1.02$ on the holdout, and $2.07$ versus $0.82$ on the
extended sample. The maximum drawdown is roughly four times
smaller on each window.

\subsection{Yield under the slippage overlay}

We described the 3-cent-per-leg slippage overlay in
Section~\ref{sec:slippage}. Its effect on yield is reported in
Table~\ref{tab:slippage}.

\begin{table}[htbp]
\centering
\small
\begin{tabular}{lrrr}
\toprule
window           & Yield (0\textcent) & Yield (3\textcent/leg) & Haircut \\
\midrule
in-sample 2020--2024 & 29.21\%  & 17.63\%  & $-40\%$  \\
holdout 2025--2026   & 35.55\%  & 25.05\%  & $-29\%$  \\
extended 2020--2026  & 31.08\%  & 20.85\%  & $-33\%$  \\
\bottomrule
\end{tabular}
\caption{Yield at zero slippage versus the 3-cent-per-leg post-hoc
slippage overlay ($\$6$ per spread for the $q = 1$ configuration,
$\$12$ per spread for $q = 2$; one-sided since spreads are held to
expiry). Yield is defined as
$\sum \mathrm{P\&L} / \sum \mathrm{wagered}$, with the wagered
denominator including the slippage cost as part of dollar at
risk.}
\label{tab:slippage}
\end{table}

The fractional haircut is large, but the strategy stays profitable
in every window even at the 3-cent-per-leg overlay. We treat
3 cents per leg as an upper bound for an investor working a
limit-order ladder on a normally-quoted SP100 chain; realised
slippage will typically be less.


\section{The \texorpdfstring{$k$}{k}-sweep}\label{sec:ksweep}

The ambiguity-aversion parameter $k$ is the one tuning knob in
the canonical intrinsic GROUND ratio. We selected $k$ on the in-sample
2020--2024 window only, by running the full strategy at a range
of $k$ values and reporting the Sharpe ratio, drawdown, and
yield for each. Once $k$ was selected, it was frozen before any
2025+ data was scored. We selected $k$ under top-$N=5$ rather than under
the threshold rule to keep the trade count constant across the sweep,
isolating the effect of $k$ on per-trade ranker quality from any
trade-count-related Sharpe effects.

\begin{table}[htbp]
\centering
\small
\begin{tabular}{rrrrrr}
\toprule
$k$ & Trades & Final & Sharpe & Max DD & Yield \\
\midrule
0.5  & 1{,}216 & \$24{,}613 & 1.58 & $-7.2\%$   & 29.2\% \\
1    & 1{,}208 & \$25{,}056 & 1.65 & $-7.2\%$   & 30.1\% \\
2    & 1{,}197 & \$24{,}864 & 1.64 & $-7.6\%$   & 29.7\% \\
5    & 1{,}141 & \$25{,}313 & 1.72 & $-7.0\%$   & 30.6\% \\
10   & 1{,}009 & \$24{,}927 & 1.75 & $-9.3\%$   & 29.9\% \\
\textbf{20 (canonical)} & \textbf{659} & \textbf{\$22{,}263} & \textbf{1.87} & $\mathbf{-6.8\%}$ & \textbf{24.5\%} \\
50   & 161    & \$14{,}690 & 1.44 & $-5.0\%$   & 9.4\% \\
100  & 52     & \$11{,}454 & 0.69 & $-3.6\%$   & 2.9\% \\
\bottomrule
\end{tabular}
\caption{The in-sample $k$-sweep at $q = 1$, under the canonical
$\GROUND = E \cdot \exp(-k D_{\mathrm{KL}})$ ranker without the
$b$-credit cap. The Sharpe ratio peaks at $k = 20$ (1.87), increasing
monotonically from $k = 0.5$ (1.58) through $k = 20$ and decreasing
monotonically through $k = 100$ (0.69). Trade counts vary as the
entropic regularization parameter $k$ increases, removing lower-confidence
intrinsic GROUND candidates while raising per-trade ranker quality. We freeze
$k = 20$ as the canonical value before any 2025+ data is scored.}
\label{tab:k-sweep}
\end{table}

A few things are worth flagging about Table~\ref{tab:k-sweep}.
First, the sweep brackets the in-sample Sharpe peak. The Sharpe
rises monotonically from $1.58$ at $k = 0.5$ through $1.87$ at $k
= 20$, then falls monotonically through $1.44$ at $k = 50$ and
$0.69$ at $k = 100$. We extended the
sweep beyond $k = 20$ for exactly this reason: the falling half
rules out an alternative explanation, that the entropic
discount is doing all the work and the Kelly EV factor is
contributing nothing. The
strategy needs both ingredients in the canonical product; pushing
the entropy penalty too hard collapses performance, just as the
single-factor baseline comparison of
Section~\ref{sec:comparison} shows that dropping the entropy
penalty entirely also degrades performance.

Second, the strategy is robust to $k$ within roughly a factor of
ten around the peak: Sharpe stays in $[1.58, 1.87]$ across $k
\in [0.5, 20]$ and drawdown improves monotonically through that
range. Outside that range the strategy degrades quickly: at $k =
100$ the Sharpe falls to $0.69$ and drawdown widens to $-3.6\%$.
The relevant consequence of this robustness for the
selection-bias concern is that with the eight grid points
$\{0.5, 1, 2, 5, 10, 20, 50, 100\}$, the in-sample Sharpe
profile identifies the peak only to within a factor of roughly
two: $k = 20$ is statistically indistinguishable from values in
$[5, 50]$ within the Sharpe noise. The in-sample peak should not
be read as a tight estimate of an optimal $k$; it is the highest
point on a slowly-varying plateau. Table~\ref{tab:k-mertens} reports
Mertens 95\% confidence intervals on the in-sample Sharpe at
$k \in \{5, 10, 20, 50\}$. All four intervals overlap on the
substantial subset $[0.74, 2.29]$, which is the statistical
content of the plateau claim.

\begin{table}[htbp]
\centering
\small
\begin{tabular}{rrrrcc}
\toprule
$k$ & Sharpe & $\gamma_3$ & $\gamma_4$ & $\sigma_{\hat S}$(Mertens) & 95\% CI \\
\midrule
$5$  & 1.72 & $+0.61$ & $+0.59$ & 0.08 & $[1.55,\,1.88]$ \\
$10$ & 1.75 & $+0.63$ & $+1.02$ & 0.09 & $[1.57,\,1.93]$ \\
$\mathbf{20}$ \textbf{(canonical)} & $\mathbf{1.87}$ & $\mathbf{+1.06}$ & $\mathbf{+3.69}$ & $\mathbf{0.12}$ & $\mathbf{[1.63,\,2.12]}$ \\
$50$ & 1.44 & $+1.42$ & $+11.00$ & 0.15 & $[1.15,\,1.73]$ \\
\bottomrule
\end{tabular}
\caption{Mertens 95\% confidence intervals on the in-sample Sharpe
at four $k$ values bracketing the in-sample peak. All four intervals
overlap on the common subset $[1.15, 1.88]$. The implication is
that $k = 20$ is not statistically distinguishable from $k \in
[5, 50]$ on the in-sample window; the choice of $k = 20$ should be
read as the highest point on a slowly-varying plateau, not as a
tight estimate of an optimal $k$. These values use threshold-based
selection matching Table~\ref{tab:k-sweep}, computed under the
uncapped configuration without the $b$-credit cap.}
\label{tab:k-mertens}
\end{table}

Third, the canonical $k = 20$ was frozen before any 2025+ data
was scored. The 2025--2026 holdout's headline performance under
the threshold rule reported in Table~\ref{tab:headline} was
achieved using a value of $k$ chosen entirely on the 2020--2024
window. A preliminary $k$-sweep was run with top-$N = 5$ selection
to identify the $k$ that maximised Sharpe; the canonical threshold
rule of $\GROUND \geq 0.10\%$ (Section~\ref{sec:threshold}) was
selected after $k$ was already frozen and operates on the same
$\GROUND$ scores produced under $k = 20$. Table~\ref{tab:k-mertens}
reports Mertens intervals under threshold-based selection for
consistency with Table~\ref{tab:k-sweep}.


\section{Single-factor baseline comparison}\label{sec:comparison}

The most important empirical question in this paper is whether the
\emph{joint} growth-and-divergence signal does work that neither
component does in isolation. To answer this question, we replace
the intrinsic GROUND ranker with each of its two factors and
re-run the strategy on the same in-sample 2020--2024 window at
$q = 1$.

\paragraph{Selection convention for this section.}
For the fair across-ranker comparison reported in Table~\ref{tab:comparison},
we use a common top-$N = 5$ per-week selection rule for all three
rankers. The $\GROUND$ ranking is applied under top-$N = 5$
selection in this section only, to enable the matched control against
which the single-factor baselines are compared.

\begin{table}[htbp]
\centering
\small
\begin{tabular}{lrrrrr}
\toprule
ranker (top-$N = 5$ selection) & Trades & Final & Sharpe & Max DD & Yield \\
\midrule
\textbf{$\GROUND$ ($k = 20$)}     & 1{,}286 & \$25{,}475 & \textbf{1.71} & $-6.8\%$  & \textbf{16.61\%} \\
$E$ alone (Kelly EV)              & 1{,}286 & \$23{,}747 & 1.43          & $-9.1\%$ & 14.64\% \\
$\KL$ alone (least-concentrated)  & 1{,}286 & \$18{,}678 & 1.00          & $-13.1\%$ & 8.45\%  \\
\bottomrule
\end{tabular}
\caption{Single-factor baseline rankers, in-sample 2020--2024,
$q = 1$, all under the canonical per-spread Kelly payoffs $\{+b,
+\alpha(b) b, -1\}$ and a common top-$N = 5$ per-week selection
rule. Under this matched convention the intrinsic
$\GROUND$ ranker outperforms both single-factor baselines on
every reported metric.}
\label{tab:comparison}
\end{table}

Under the common top-$N = 5$ convention, the joint ranker beats
the $E$-alone baseline on the Sharpe ratio ($1.71$ versus
$1.43$), on yield ($16.61\%$ versus $14.64\%$), and on drawdown
($-6.8\%$ versus $-9.1\%$). It beats the $\KL$-alone baseline by
a wider margin on every metric.

The $E$-alone row is the more informative comparison, because it
isolates the contribution of the entropic discount. Without the
$\exp(-k \KL)$ factor, the strategy loses control of the tails:
drawdown deepens by $34\%$, from $-6.8\%$ to $-9.1\%$, and
Sharpe drops by $0.28$. This is consistent with the motivation in
Section~\ref{sec:intro}. The entropy penalty discounts trades whose
Kelly-EV advantage depends on sharp implied-distribution
concentration, which is exactly the kind of edge that does not
survive estimation noise. The $\KL$-alone row collapses on yield,
which is also expected: ranking by lowest divergence and ignoring
growth gives the most-uniform distribution, which carries no
particular reason to be profitable.

The two component signals are individually weak. Their
multiplicative combination via $\GROUND = E \cdot \exp(-k \KL)$ is
where the empirical edge lives.

\section{Threshold-based selection: user-controlled risk-adjusted return}\label{sec:threshold}

The headline results of Section~\ref{sec:results} use the canonical
selection rule $\GROUND \geq \tau$ with $\tau = 0.10\%$. This
section motivates that choice.

The threshold rule has a direct economic reading. Under the
canonical Kelly-EV-with-entropic-discount form, $\GROUND$ is
dimensionless and reads as the variance-adjusted expected return per trade
after the entropic ambiguity discount (Section~\ref{sec:framework}). The cutoff
$\GROUND \geq \tau$ is therefore equivalent to imposing a per-trade
hurdle of $\tau$. An investor with a $0.10\%$ hurdle rate gets
a strategy that trades exactly the candidates clearing that
hurdle; an investor with a higher hurdle gets fewer,
higher-quality trades. The control surface is the hurdle rate
itself.

To select $\tau$ we ran the strategy under $\GROUND \geq \tau$ for
$\tau \in \{0,\,0.05\%,\,0.10\%,\,0.25\%,\,0.50\%,\,1\%,\,2\%,\,5\%\}$
on the 2020--2024 in-sample window at $q = 1$ with all other
configuration canonical (Greek filters, regime gate, $k = 20$).

\begin{table}[htbp]
\centering
\small
\begin{tabular}{rrrrrrr}
\toprule
$\tau$ & Trades & t/wk & Final & Sharpe & Max DD & Yield \\
\midrule
$0.00\%$ & 4{,}731 & 18.27 & \$32{,}603 & 0.70 & $-29.9\%$ & 5.59\%  \\
$0.05\%$ & 1{,}082 &  4.18 & \$26{,}453 & 1.66 & $-10.6\%$ & 21.02\% \\
$\mathbf{0.10\%}$ \textbf{(canonical)} & \textbf{731} & \textbf{2.81} & \textbf{\$24{,}653} & \textbf{1.85} & $\mathbf{-7.7\%}$ & \textbf{29.21\%} \\
$0.25\%$ &   373    &  1.43 & \$19{,}232 & 1.97 & $-6.6\%$  & 36.24\% \\
$0.50\%$ &   197    &  0.76 & \$16{,}537 & 1.70 & $-6.2\%$  & 50.78\% \\
$1.00\%$ &   100    &  0.39 & \$14{,}511 & 1.70 & $-2.5\%$  & 70.58\% \\
$2.00\%$ &    37    &  0.14 & \$12{,}139 & 1.29 & $-1.4\%$  & 95.46\% \\
$5.00\%$ &    10    &  0.04 & \$10{,}334 & 0.49 & $-1.4\%$  & 81.40\% \\
\bottomrule
\end{tabular}
\caption{In-sample 2020--2024 threshold sweep at $q = 1$. Each row
shows results for a different $\GROUND$ cutoff $\tau$, with columns
reporting the number of selected candidates, trades per week, final
equity, Sharpe, max drawdown, and yield.}
\label{tab:threshold-sweep}
\end{table}

The Sharpe profile rises sharply from $\tau = 0$ as low-quality
candidates are filtered out, peaks at approximately
$\tau = 0.25\%$ (in-sample Sharpe $1.97$), and falls past that
point as the strategy runs out of trades. The $\tau = 0$ row is
informative: ranking with no cutoff at all produces a Sharpe of
$0.70$ with a $-30\%$ drawdown, which establishes that the bottom
of the $\GROUND$ distribution is not informative and that some
selectivity is essential.

We adopt $\tau = 0.10\%$ rather than the Sharpe-maximizing
$\tau = 0.25\%$, and we want to be explicit about the criterion.
The choice was not Sharpe-max on the in-sample window; it was a
trade-count floor. At $\tau = 0.25\%$ the strategy retains $1.43$
trades per week, which we judged too thin to support the
$\sqrt{N}$ diversification structure of the weekly Sharpe
estimator. At $\tau = 0.10\%$ the strategy keeps approximately
$2.81$ trades per week with in-sample Sharpe $1.85$, drawdown
$-7.7\%$, and yield $29.21\%$. The buffer below the in-sample
Sharpe peak is the conservative side of the selection.

The in-sample Sharpe of $1.85$ at $\tau = 0.10\%$ is also higher
than the $1.71$ Sharpe of the intrinsic GROUND ranker under top-$N = 5$
selection at the same $k = 20$ reported in Table~\ref{tab:comparison}. The mechanism is
direct: the threshold rule filters out marginal top-$5$ picks
whose $\GROUND$ scores fall below $0.10\%$, which are
disproportionately the lower-quality candidates in each week's
top-$5$. Removing them concentrates the strategy on the part of
the candidate distribution where the ranker is genuinely
informative (Section~\ref{sec:decile}).


\section{Ranker quality across the full candidate distribution}\label{sec:decile}

The headline results of Section~\ref{sec:results} measure the
strategy with the canonical $\GROUND \geq 0.10\%$ threshold
applied, which is the configuration the strategy actually trades.
In this section we ask a different question: how informative is
$\GROUND$ as a \emph{per-candidate} signal, across the full pool
of all candidates that survive the Greek, regime, and
open-interest filters?

To answer this, we re-run with $\texttt{top\_n} = \infty$, that
is, with no per-week pick limit. The candidate pool that emerges is
$6{,}130$ candidate trades over the 2020--2026 period. This is the
largest sample we can produce for the diagnostic, and the one in
which the ranker's per-candidate predictive content is most
directly visible.

\subsection{Pooled deciles}

\begin{table}[htbp]
\centering
\small
\begin{tabular}{rrrrr}
\toprule
$\GROUND$ decile & mean yield (\%) & win rate (\%) & $n$ & mean $\GROUND$ (\%) \\
\midrule
1 (lowest)  & $+2.41$  & 50.5 & 727 & $0.00\%$ \\
2           & $-3.57$  & 45.8 & 561 & $0.00\%$ \\
3           & $+1.69$  & 48.7 & 567 & $0.00\%$ \\
4           & $+5.54$  & 50.2 & 606 & $0.01\%$ \\
5           & $-0.83$  & 46.9 & 612 & $0.01\%$ \\
6           & $+3.98$  & 48.3 & 607 & $0.02\%$ \\
7           & $+10.89$ & 50.1 & 613 & $0.05\%$ \\
8           & $+14.53$ & 48.9 & 611 & $0.12\%$ \\
9           & $+18.65$ & 47.3 & 613 & $0.12\%$ \\
\textbf{10 (highest)}& $\mathbf{+76.71}$ & \textbf{44.7} & \textbf{613} & $\mathbf{1.33\%}$ \\
\bottomrule
\end{tabular}
\caption{Per-trade yield, win rate, and candidate count by
$\GROUND$ decile across the full candidate distribution. All SP100
vertical credit spreads passing the canonical filters over
2020--2026 are included ($n = 6{,}130$ candidates, $q = 1$). The
Spearman rank correlation between $\GROUND$ and yield across all
candidates is $+0.118$.}
\label{tab:decile}
\end{table}

The pooled Spearman correlation between $\GROUND$ and per-trade
yield is $+0.118$. This is modest. We acknowledge it directly:
across all candidates, the per-candidate predictive content is
not large. The picture in Table~\ref{tab:decile} is clearer than
the pooled correlation. The signal concentrates in the top
decile, which beats the bottom decile by roughly $32 \times$ on
mean yield ($+76.71\%$ versus $+2.41\%$). Win rates are flat
across deciles, in the range $44\%$ to $51\%$. The ranker is not
selecting candidates that win more often; it is selecting
candidates that win bigger when they do win, and lose smaller
when they lose.

\subsection{Sub-period robustness}

A natural concern about Table~\ref{tab:decile} is that the
top-decile edge could be a 2020--2024 artefact that fails to
reproduce on the strict 2025--2026 holdout. We address this
concern directly in Table~\ref{tab:decile-split} by splitting the
deciles by sub-period.

\begin{table}[htbp]
\centering
\small
\begin{tabular}{r rr rr}
\toprule
 & \multicolumn{2}{c}{in-sample 2020--2024} & \multicolumn{2}{c}{holdout 2025--2026} \\
$\GROUND$ decile & mean yield (\%) & $n$ & mean yield (\%) & $n$ \\
\midrule
1 (lowest)  & $+4.12$  & 550 & $-0.16$  & 154 \\
2           & $-3.01$  & 429 & $-9.34$  & 126 \\
3           & $-6.75$  & 475 & $+18.86$ & 131 \\
4           & $+6.78$  & 453 & $+12.99$ & 138 \\
5           & $-1.84$  & 477 & $+2.67$  & 135 \\
6           & $+0.79$  & 479 & $+13.15$ & 137 \\
7           & $+8.61$  & 475 & $+22.49$ & 134 \\
8           & $+14.52$ & 474 & $+19.86$ & 138 \\
9           & $+15.63$ & 476 & $+27.00$ & 136 \\
\textbf{10 (highest)} & $\mathbf{+62.96}$ & \textbf{477} & $\mathbf{+121.70}$ & \textbf{136} \\
\midrule
\multicolumn{1}{l}{Spearman $\rho$} & \multicolumn{2}{c}{$+0.108$} & \multicolumn{2}{c}{$+0.153$} \\
\bottomrule
\end{tabular}
\caption{Decile mean yield split by sub-period: in-sample 2020--2024 ($n = 4{,}765$ candidates) and strict 2025--2026 holdout ($n = 1{,}365$ candidates).}
\label{tab:decile-split}
\end{table}

The split confirms the pooled picture and strengthens it on the
out-of-sample side. The top-decile yield is substantially higher in the holdout
than in-sample ($+121.70\%$ versus $+62.96\%$), the bottom-decile
yield is worse in the holdout ($-0.16\%$ versus $+4.12\%$), and
the Spearman correlation is stronger in the holdout
($+0.153$ versus $+0.108$). The top-decile-edge result is
therefore robust to the sub-period split.

The mid-decile rows (deciles $3$, $5$, and $6$) exhibit large
sign-flips between the two sub-periods. This is the expected
behaviour of a noisy mid-of-distribution signal in a 70-week
holdout, and it is consistent with the reading that the ranker
concentrates its predictive content at the extremes of the
distribution. The instability in the middle of the table is
diagnostic noise, not strategy risk: under the canonical
threshold rule the strategy only selects from deciles whose
$\GROUND$ clears $\tau = 0.10\%$, which in practice means decile
$10$ and the top of decile $9$. The mid-deciles never enter the
P\&L.

\subsection{Interpretation: the decile pattern and entropic regularization}

The most distinctive empirical feature of
Table~\ref{tab:decile} is that win rate and yield run in opposite directions across deciles. The top decile has the worst win rate at $44.7\%$, while decile 7 reaches $50.1\%$; yet the top decile's mean yield of
$+76.71\%$ dwarfs decile 7's $+10.89\%$ and beats the bottom by
approximately $32\times$. This inversion (lowest win frequency, highest payoff magnitude) is the empirical signature of
payoff-asymmetry selection, not probability-timing selection. The
$\GROUND$ ranker does not select candidates that win more often; it
selects candidates that win bigger when they do win and lose
smaller when they lose.

This is the empirical signature one would expect from the
ambiguity-aversion regularization that motivates the criterion in
Section~\ref{sec:intro}. The structure of $\GROUND = E \cdot
\exp(-k \cdot \KL)$ selects for robustness: a candidate with a sharply concentrated
implied distribution can have high Kelly EV $E$, but the entropic
discount $\exp(-k \cdot \KL)$ pushes it down the ranking. The top
$\GROUND$ candidates are therefore not the most-concentrated bets
but candidates whose Kelly EV is positive while their implied
distribution sits close enough to the uniform reference that the
ambiguity discount remains modest. These are bets whose payoff
profile is informationally well-conditioned: when they pay off,
they pay off in proportion to their Kelly EV; when they don't,
their losses are not amplified by the sharp probability
concentration that would have made them tail-fragile in the first
place.

Theorem~\ref{thm:memm} shows that under Kelly-optimal sizing, in the small-edge asymptotic limit, the
reservation-price approximation is asymptotically proportional to expected payoff under any belief, with proportionality constant $\tfrac{1}{2}$.
The intrinsic GROUND ranker with $k > 0$ applies an entropy penalty that further discounts
prices in proportion to distributional divergence from uniform. This is a
robustness mechanism, not a MEMM-specific construction: it applies identically
to any belief measure, whether MEMM, risk-neutral, or estimated empirical.
The empirical decile pattern (flat win rates, magnitude-asymmetric payoffs)
reflects the intrinsic GROUND ranker's joint optimization of edge and calibration, not a special
relationship to MEMM. The regularization works the same for any choice of belief measure.


\section{Discussion}\label{sec:discussion}

\subsection{Limitations}

We acknowledge three limitations of the present study.

\paragraph{Holdout length.}
The holdout is 70 weekly observations, which is a short
statistical record. The Mertens 95\% confidence interval on the
holdout Sharpe is wide at $[1.39, 4.21]$ (Table~\ref{tab:sharpe-ci}). The extended-sample row
is the more defensible external number; the holdout's role here is
to confirm that $k = 20$, frozen at in-sample selection, did not
collapse on the strict out-of-sample window. It did not.

\paragraph{Slippage upper bound.}
The 3-cent-per-leg overlay is a conservative upper bound on
per-leg execution slippage. Realized slippage on normally-quoted
SP100 chains, worked through a limit-order ladder, will typically
be less. Yields at the 3-cent-per-leg bound are $17.63\%$
in-sample and $25.05\%$ on the holdout
(Table~\ref{tab:slippage}).

\paragraph{Single regime indicator.}
The strategy uses one binary macro filter, namely SPY closing
above its 100-day SMA. Per-ticker regime gates,
VIX-conditioned filters, and gap filters were evaluated during
development and were found to be either redundant with the SPY
gate or return-degrading. The canonical configuration accordingly
uses the single SPY filter only.

\subsection{Open problems}

Three directions in which the framework of this paper could be
extended:

\paragraph{Multi-period extension.}
The single-period myopic treatment matches weekly practice. Under
bankroll-coupled fractional-Kelly sizing, where the next period's
wager scales with current bankroll, the multi-period extension is
the standard Bellman recursion
\begin{equation*}
V_t(b) \;=\; \max_{\C \in \mathcal{M}_t} \Bigl\{\, J_k(p_\C, x_\C)
   \;+\; \beta \cdot \E_{p_\C}\!\bigl[ V_{t+1}\!\bigl(b \cdot M(p_\C, x_\C; w^\star)\bigr) \bigr] \,\Bigr\},
\end{equation*}
with terminal condition $V_T(b) = 0$ and per-period discount
$\beta \in (0, 1]$. The selection at each step is the
$\GROUND$-style argmax plus a continuation term that depends on
the realised wealth multiplier. Under the bankroll-independent
flat sizing of the canonical configuration, the continuation
term is constant in the choice $\C$ and the per-period myopic
rule already used here is exactly optimal. The
fractional-Kelly variant where bankroll matters is the natural
candidate for the dynamic-programming treatment, and an empirical
study of that variant under realistic bankroll-growth paths
would establish optimality in a stronger sense.

\paragraph{Continuous-state limit.}
The canonical form $\GROUND = E \cdot \exp(-k \KL)$ extends to
continuous outcome spaces with the differential KL replacing the
discrete KL and the Kelly first-order condition becoming an
integral equation. The tight integrability conditions for
unbounded payoffs (long straddles, naked tail wings) remain open.

\paragraph{Alternative reference measures.}
The uniform reference $u_n$ is used throughout this paper.
Alternative references, for example the risk-neutral measure
derived from the most-liquid at-the-money contract on the same
underlying, are economically defensible and would recover the
comparative GROUND form of Mercurio et al.\ \cite{mercurio2020binary}. The cost of
moving away from uniform is the loss of the scale-anchored
interpretation of $\KL \in [0, \ln n]$.


\appendix

\section{Deferred proofs}\label{app:proofs}

\subsection{Existence and uniqueness of \texorpdfstring{$w^\star$}{w-star}}

\begin{lemma}[Properties of $w^\star$]\label{lem:wstar}
Let $\C = (x, p, c)$ be a contingent claim with $\min_\omega
x(\omega) = -1$, $\max_\omega x(\omega) = +1$, and $p(\omega) > 0$
for all $\omega$ with $x(\omega) \neq 0$. Then the expected
log-growth $\ell(\cdot; p, x): [0, 1) \to \R$ defined in
\eqref{eq:loggrowth} is strictly concave and admits a unique
maximizer $w^\star \in [0, 1)$ if and only if $\E_p[x] > 0$, in
which case $w^\star \in (0, 1)$.
\end{lemma}

\begin{proof}
The function $w \mapsto \log(1 + w x(\omega))$ is concave for any
$x(\omega) > -1$, and strictly concave when $x(\omega) \neq 0$. The
sum $\ell$ is therefore concave on $[0, 1)$, and strictly concave
under the support condition. The first-order condition
$\ell'(w; p, x) = \sum_\omega p(\omega) x(\omega) / (1 + w
x(\omega)) = 0$ has a unique solution in $(0, 1)$ if and only if
$\ell'(0) = \E_p[x] > 0$ and $\ell'(w) \to -\infty$ as $w \to 1^-$,
which holds whenever $p(\omega) > 0$ for some $\omega$ with
$x(\omega) = -1$. Strict concavity then gives uniqueness of the
interior optimizer. When $\E_p[x] \le 0$, strict concavity implies
that $\ell$ is non-increasing on $[0, 1)$ and $w^\star = 0$
optimizes. \qed
\end{proof}

\subsection{Closed-form \texorpdfstring{$w^\star$}{w-star} in the three-outcome case}

For the three-outcome credit-spread setting with payoffs
$\{+b, +\alpha b, -1\}$ and probabilities $\{p, r_o, q\}$, the
first-order condition becomes a quadratic in $w$ of the form $A w^2
+ B w + C = 0$, with
\begin{align*}
A &= -\alpha b^2, \\
B &= \alpha b^2 (p + r_o) - b \bigl[ p + r_o \alpha + q(1 + \alpha) \bigr], \\
C &= p b + r_o \alpha b - q.
\end{align*}
The unique root in $(0, 1)$ is
\begin{equation*}
w^\star \;=\; \frac{-B - \sqrt{B^2 - 4AC}}{2A}.
\end{equation*}

\subsection{Proof of Propositions~\ref{prop:kelly-limit}~and~\ref{prop:maxent-limit}}

\paragraph{Kelly limit.} $\GROUND(p, x, k) = E(p, x) \cdot \exp(-k
\KL(p \| u_n)) \to E(p, x)$ as $k \to 0$, uniformly in $(p, x)$
on a finite menu. The arg-max of a uniformly convergent sequence
on a finite set converges to the arg-max of the limit, modulo
ties.

\paragraph{Maximum-entropy limit.} Take natural logarithms: $\ln
\GROUND = g - k \KL = J_k$, so for large $k$ the $-k \KL$ term
dominates the bounded $g$ on a finite menu, and $\arg\max_i J_k
(p_i, x_i) \to \arg\min_i \KL(p_i \| u_n)$. Strict monotonicity of
the exponential preserves the arg-max under the multiplicative
form.


\section{Materials and Methods}\label{sec:materials}

The option chain data used in the backtest is end-of-day SP100
single-name option pricing data sourced from a commercial historical
options-data vendor, covering the period 2020-01-01 through
2026-05-04. The fields used in this study are the per-leg bid, ask,
mid-price, open interest, implied volatility, and the standard
Black-Scholes Greeks. SPY daily close prices, used to construct the
$100$-day simple moving average regime indicator, are sourced from a
standard dividend-adjusted public daily-price feed.

All quantities in this paper are derived from these two sources. The
backtest, the in-sample $k$-sweep, the single-factor baseline
comparison, and the decile diagnostic are produced by a Python
implementation maintained by the author.

\paragraph{Author Contributions:}
Conceptualization, P.J.M.; Data curation, P.J.M.; Formal analysis,
P.J.M.; Investigation, P.J.M.; Methodology, P.J.M.; Resources,
P.J.M.; Software, P.J.M.; Validation, P.J.M.; Visualization, P.J.M.;
Writing---original draft, P.J.M.; Writing---review \& editing,
P.J.M. The author has read and agreed to the published version of
the manuscript.

\paragraph{Funding:}
This research received no external funding.

\paragraph{Conflicts of Interest:}
The author declares no conflict of interest.

\paragraph{Abbreviations:}
The following abbreviations are used in this manuscript:

\medskip
\begin{tabbing}
\hspace{2.5cm}\=\kill
CI       \> Confidence interval \\
DKL      \> Kullback-Leibler divergence (KL) \\
DTE      \> Days to expiry \\
EV       \> Kelly EV (variance-adjusted expected return per trade) \\
FOC      \> First-order condition \\
GROUND   \> Growth rate over uniform divergence \\
ITM      \> In-the-money \\
KKT      \> Karush-Kuhn-Tucker (optimality conditions) \\
KL       \> Kullback-Leibler \\
MEMM     \> Minimum-entropy martingale measure \\
OI       \> Open interest \\
P\&L     \> Profit and loss \\
ROI      \> Return on investment \\
SMA      \> Simple moving average \\
SP100    \> Standard \& Poor's 100 (equity index) \\
SPY      \> SPDR S\&P 500 ETF Trust (ticker symbol) \\
VIX      \> CBOE Volatility Index \\
\end{tabbing}


\bibliographystyle{unsrtnat}
\begin{thebibliography}{99}

\bibitem{kelly1956}
Kelly, J. A New Interpretation of Information Rate.
\emph{Bell Syst.\ Tech.\ J.} \textbf{1956}, \emph{35}, 917--926.
\href{https://doi.org/10.1002/j.1538-7305.1956.tb03809.x}{[CrossRef]}

\bibitem{breeden1978}
Breeden, D.T.; Litzenberger, R.H. Prices of State-Contingent Claims
Implicit in Option Prices. \emph{J.\ Bus.} \textbf{1978}, \emph{51},
621--651.
\href{https://doi.org/10.1086/296025}{[CrossRef]}

\bibitem{maclean2011}
MacLean, L.C.; Thorp, E.O.; Ziemba, W.T., Eds. \emph{The Kelly
Capital Growth Investment Criterion}; World Scientific: Singapore,
\textbf{2011}.

\bibitem{mercurio2020binary}
Mercurio, P.J.; Wu, Y.; Xie, H. Portfolio Optimization for Binary
Options Based on Relative Entropy. \emph{Entropy} \textbf{2020},
\emph{22}, 752.
\href{https://doi.org/10.3390/e22070752}{[CrossRef]}

\bibitem{frittelli2000}
Frittelli, M. The Minimal Entropy Martingale Measure and the
Valuation Problem in Incomplete Markets. \emph{Math.\ Financ.}
\textbf{2000}, \emph{10}, 39--52.
\href{https://doi.org/10.1111/1467-9965.00079}{[CrossRef]}

\bibitem{mercurio2020portfolio}
Mercurio, P.J.; Wu, Y.; Xie, H. An Entropy-Based Approach to
Portfolio Optimization. \emph{Entropy} \textbf{2020}, \emph{22},
332.
\href{https://doi.org/10.3390/e22030332}{[CrossRef]}

\bibitem{follmer2002}
F\"ollmer, H.; Schweizer, M. Convex Measures of Risk and Trading
Constraints. \emph{Finance Stoch.} \textbf{2002}, \emph{6},
429--447.
\href{https://doi.org/10.1007/s007800200072}{[CrossRef]}

\bibitem{mercurio2020}
Mercurio, P.J.; Wu, Y.; Xie, H. Option Portfolio Selection with
Generalized Entropic Portfolio Optimization. \emph{Entropy}
\textbf{2020}, \emph{22}, 805.
\href{https://doi.org/10.3390/e22080805}{[CrossRef]}

\bibitem{hansen2008robustness}
Hansen, L.P.; Sargent, T.J. \emph{Robustness}; Princeton University
Press: Princeton, NJ, USA, \textbf{2008}.

\bibitem{maccheroni2006ambiguity}
Maccheroni, F.; Marinacci, M.; Rustichini, A. Ambiguity Aversion,
Robustness, and the Variational Representation of Preferences.
\emph{Econometrica} \textbf{2006}, \emph{74}, 1447--1498.
\href{https://doi.org/10.1111/j.1468-0262.2006.00716.x}{[CrossRef]}

\bibitem{delbaen2002}
Delbaen, F.; Grandits, P.; Rheinl\"ander, T.; Samperi, D.;
Schweizer, M.; Stricker, C. Exponential Hedging and Entropic
Penalties. \emph{Math.\ Financ.} \textbf{2002}, \emph{12},
99--123.
\href{https://doi.org/10.1111/1467-9965.02001}{[CrossRef]}

\bibitem{fujiwara2003}
Fujiwara, T.; Miyahara, Y. The Minimal Entropy Martingale Measures
for Geometric L\'evy Processes. \emph{Finance Stoch.} \textbf{2003},
\emph{7}, 509--531.
\href{https://doi.org/10.1007/s007800200099}{[CrossRef]}

\bibitem{lo2002sharpe}
Lo, A.W. The Statistics of Sharpe Ratios. \emph{Financ.\ Anal.\ J.}
\textbf{2002}, \emph{58}, 36--52.
\href{https://doi.org/10.2469/faj.v58.n4.2453}{[CrossRef]}

\bibitem{mertens2002sharpe}
Mertens, E. Comments on Variance of the IID Estimator in Lo (2002).
University of Basel, \textbf{2002}.

\end{thebibliography}

\end{document}
